% Session 1 of "From Paper to Proof: Formalizing Mathematics in Lean with AI".
% Build: Overleaf with the pdfLaTeX compiler (or pdflatex, twice)
\documentclass[10pt, aspectratio=169]{beamer}
\usepackage{bgsmathlean}
\leanfile{Session1.lean}

\title{From Paper to Proof}
\subtitle{Session 1: Lean and Mathlib}
\author{Michail Karatarakis}
\institute{BGSMath Graduate Course \textperiodcentered{} UPC, Barcelona}
\date{10 November 2026}

\begin{document}

\maketitle

%============================================================================
\begin{frame}[fragile]{The course}
\sessionplan{1}

\vspace{0.4em}
Each session: a lecture with live demonstrations in Lean. You redo them, and the exercises, at home.
\end{frame}

%----------------------------------------------------------------------------
\begin{frame}[fragile]{What formalizing means}
\begin{itemize}
\item Writing definitions, statements and proofs in a language a computer can read.
\item Every step of the proof is checked by a program, Lean, down to the axioms.
\item Building on a shared library, Mathlib: the basic definitions and theorems are already there.
\item The result is a statement, and a proof of it that Lean accepts.
\end{itemize}

\vspace{0.8em}
\begin{keypoint}
\normalsize A formal proof is a proof that a computer has checked, line by line.
\end{keypoint}
\end{frame}

%----------------------------------------------------------------------------
\begin{frame}[fragile]{What formalizing does not mean}
\begin{itemize}
\item \textbf{Not that the theorem is the one you meant.} Lean checks the proof against the statement as written. Whether the statement says what the paper says is for you to check. This is today's main topic.
\item \textbf{Not that the computer finds the proofs.} Automation and AI can propose proofs; Lean only checks them.
\item \textbf{Not a replacement for the paper.} A formal proof shows that something is true, not why. The ideas and the explanation stay on paper.
\item \textbf{Not programming.} This course assumes no Lean and no programming. What it needs is careful reading.
\item \textbf{Not all or nothing.} You can state a theorem without proving it, or prove it assuming a step you cannot do yet.
\end{itemize}
\end{frame}

%----------------------------------------------------------------------------
\begin{frame}[fragile]{Why now}
In the last twelve months, AI has written a large part of the Lean code in several major formalizations.

\vspace{0.3em}
\small
\renewcommand{\arraystretch}{1.25}
\begin{tabularx}{\linewidth}{@{}>{\raggedright\arraybackslash\bfseries}p{0.2\linewidth}X@{}}
\toprule
Fermat's Last Theorem & A complete Lean proof, produced by Claude and drawing on Kevin Buzzard's FLT project.\\
Sphere packing & Viazovska's theorems in dimensions 8 and 24, completed with Math Inc's Gauss from a project started by Hariharan and Viazovska. Gauss's part is disputed (session 3).\\
Unit distances & An AI-found disproof of Erdős's unit distance conjecture, formalized in full by Boris Alexeev with ChatGPT.\\
Navier--Stokes & OpenAI announced Navier--Stokes blowup with forcing, in Lean. Whether it matches the Millennium Prize statement needs a human check (Hales).\\
\bottomrule
\end{tabularx}

\source{Anthropic (Sep 2026); Math, Inc.\ (Feb 2026); K.\ Buzzard, Xena blog (Jul 2026); OpenAI (Sep 2026).}
\end{frame}

%----------------------------------------------------------------------------
\begin{frame}[fragile]{The one idea of today}
\begin{keypoint}
Lean checks the \textbf{proof}. You check the \textbf{statement}.
\end{keypoint}

\vspace{0.8em}
\begin{itemize}
\item An AI can hand you a proof that Lean accepts.
\item Whether it proves the theorem \emph{you meant} is for you to decide.
\item So the most useful skill is not writing Lean but \textbf{reading} it.
\end{itemize}

\vspace{0.6em}
\begin{paper}[T.\ Hales, on Terence Tao's blog, 9 October 2026]
\itshape ``A proof in Lean should not be accepted until a human audit is performed to ensure statement fidelity. [\dots] This task is generally massively easier than checking the proof itself.''
\end{paper}
\end{frame}

%----------------------------------------------------------------------------
\begin{frame}[fragile]{The game}
\begin{center}
\begin{minipage}{0.65\textwidth}
\lean{flt}
\end{minipage}
\end{center}

\begin{itemize}
\item Fermat's Last Theorem, in Lean.
\item \lc{sorry} means: proof missing.
\item Read the statement. Why the hypotheses \lc{a ≠ 0}? Without them, $0^n + b^n = b^n$.
\end{itemize}

\vspace{0.4em}
\begin{keypoint}
\normalsize The game: replace each \lc{sorry} by a proof Lean accepts. Today: the rules.
\end{keypoint}
\end{frame}

%============================================================================
\section{What Lean is}

\begin{frame}[fragile]{A proof assistant}
\begin{columns}[c]
\begin{column}{0.52\textwidth}
\begin{itemize}
\item You write definitions, statements and proofs in a formal language.
\item Lean checks every step. A small \textbf{kernel} has the final word.
\item Tactics, automation and AI may be as clever or as wrong as they like: whatever they produce is checked by the kernel.
\item Lean is also a programming language, so you can compute with your definitions.
\end{itemize}
\end{column}
\begin{column}{0.44\textwidth}
\centering
\begin{tikzpicture}[every node/.append style={text=ink}]
  \node[layer=paperrule] (you)  at (0,3.0) {Your file\\[-2pt]{\scriptsize definitions, theorems, proofs}};
  \node[layer=leanblue]  (lib)  at (0,2.0) {Mathlib\\[-2pt]{\scriptsize the mathematical library}};
  \node[layer=muted]     (elab) at (0,1.0) {Lean: elaborator and tactics\\[-2pt]{\scriptsize turn what you write into a proof term}};
  \node[layer=alertred, line width=1pt] (kern) at (0,0) {Kernel\\[-2pt]{\scriptsize checks the proof term}};
  \node[font=\scriptsize, text=alertred, anchor=west] at (kern.east) {\ trusted};
\end{tikzpicture}
\end{column}
\end{columns}
\end{frame}

%----------------------------------------------------------------------------
\begin{frame}[fragile]{Mathlib}
\begin{columns}[T]
\begin{column}{0.4\textwidth}
\begin{tikzpicture}[stat/.style={font=\bfseries\LARGE, text=leanblue, anchor=east},
                    lab/.style={font=\small, anchor=west, text=ink}]
  \node[stat] at (0,1.6) {139{,}422}; \node[lab] at (0.15,1.6) {definitions};
  \node[stat] at (0,0.8) {290{,}718}; \node[lab] at (0.15,0.8) {theorems};
  \node[stat] at (0,0)   {772};       \node[lab] at (0.15,0)   {contributors};
\end{tikzpicture}
\end{column}
\begin{column}{0.56\textwidth}
\begin{itemize}
\item One library, written by the community, reviewed by maintainers.
\item Algebra, number theory, analysis, measure theory and probability, topology, geometry, combinatorics, category theory, \dots
\item Everything is stated in the \emph{greatest generality} that is convenient: expect rings, fields and ordered structures rather than just $\mathbb{R}$.
\end{itemize}
\end{column}
\end{columns}

\source{Mathlib statistics page, October 2026.}
\end{frame}

%----------------------------------------------------------------------------
\begin{frame}[fragile]{Getting started}
\begin{columns}[T]
\begin{column}{0.31\textwidth}
\textbf{\color{leanblue}In the browser}
\begin{itemize}
\item \href{https://live.lean-lang.org}{live.lean-lang.org}
\item Mathlib included, nothing to install.
\item Enough for sessions 1 and 2.
\end{itemize}
\end{column}
\begin{column}{0.31\textwidth}
\textbf{\color{leanblue}On your laptop}
\begin{itemize}
\item VS Code with the \emph{Lean 4} extension.
\item \lc{lake new myproject math} creates a project with Mathlib.
\item For sessions 3 and 4: see the setup sheet.
\end{itemize}
\end{column}
\begin{column}{0.31\textwidth}
\textbf{\color{leanblue}The course files}
\begin{itemize}
\item The code on the slides, one Lean file per session.
\item Exercises with \lc{sorry} to fill in, at home.
\item Solutions after the session.
\end{itemize}
\end{column}
\end{columns}

\vspace{1em}
\begin{keypoint}
\normalsize The \textbf{infoview} panel next to your code shows what Lean knows at the cursor. Keep it open.
\end{keypoint}
\source{If opening a file makes Lean start building Mathlib, stop it and run \texttt{lake exe cache get} in the project folder.}
\end{frame}

%============================================================================
\section{First steps}

\begin{frame}[fragile]{First contact\livedemo}
\lean{first}

\vspace{0.3em}
\begin{itemize}
\item \lc{#eval} computes, \lc{#check} tells you the type of an expression.
\item \lc{Nat.Prime 7} is a \emph{statement}: it has type \lc{Prop}. It is not yet a proof of anything.
\end{itemize}

\vspace{0.3em}
\lean{rfl}
\begin{itemize}
\item A \emph{proof} of \lc{2 + 2 = 4}. Here \lc{rfl} says: both sides compute to the same value.
\end{itemize}
\end{frame}

%----------------------------------------------------------------------------
\begin{frame}[fragile]{Anatomy of a theorem}
\centering
\begin{tikzpicture}[node distance=0pt]
  \node[part=leanblue] (kw) {theorem};
  \node[part=okgreen, right=4pt of kw] (nm) {sq\_add\_sq\_ge};
  \node[part=paperrule, right=4pt of nm] (ar) {(a b : ℝ)};
  \node[font=\ttfamily\small, right=3pt of ar] (c1) {:};
  \node[part=alertred, right=3pt of c1] (st) {2 * a * b ≤ a \textasciicircum{} 2 + b \textasciicircum{} 2};
  \node[font=\ttfamily\small, right=3pt of st] (as) {:= by};
  \node[part=muted, anchor=north west] (pf) at ($(kw.south west)+(0.45,-0.25)$) {nlinarith [sq\_nonneg (a - b)]};

  \node[lbl=leanblue] (lkw) at ($(kw.north)+(0,1.0)$) {keyword};
  \node[lbl=okgreen] (lnm) at ($(nm.north)+(0,1.0)$) {name, to use\\it later};
  \node[lbl=ink] (lar) at ($(ar.north)+(0.2,1.0)$) {variables and\\hypotheses};
  \node[lbl=alertred] (lst) at ($(st.north)+(0,1.0)$) {the statement:\\this is what you check};
  \node[lbl=muted] (lpf) at ($(pf.south)+(0,-0.9)$) {the proof: this is what Lean checks};
  \draw[arr=leanblue] (lkw) -- (kw);
  \draw[arr=okgreen] (lnm) -- (nm);
  \draw[arr=paperrule] (lar) -- (ar);
  \draw[arr=alertred] (lst) -- (st);
  \draw[arr=muted] (lpf) -- (pf);
\end{tikzpicture}

\vspace{1em}
\raggedright
\begin{itemize}
\item \lc{example} is a \lc{theorem} without a name.
\item \lc{by} switches to \emph{tactic mode}: a sequence of instructions that build the proof step by step.
\end{itemize}
\end{frame}

%----------------------------------------------------------------------------
\begin{frame}[fragile]{From paper to Lean}
\begin{columns}[T]
\begin{column}{0.38\textwidth}
\begin{paper}
\textbf{Claim.} For all real $a, b$, \ $2ab \le a^2 + b^2$.

\textbf{Proof.} We have $(a-b)^2 \ge 0$, and
\[
\begin{aligned}
2ab &= a^2 + b^2 - (a - b)^2\\
    &\le a^2 + b^2. \qquad\square
\end{aligned}
\]
\end{paper}
\end{column}
\begin{column}{0.58\textwidth}
\leanlabel
\lean{calc}
\end{column}
\end{columns}

\vspace{0.2em}
\begin{itemize}
\item \lc{have} records an intermediate fact; \lc{calc} is a chain of (in)equalities, each with its own justification.
\item \lc{ring} proves identities in commutative rings; \lc{linarith} does linear arithmetic with the hypotheses.
\item Short version: \lc{nlinarith [sq_nonneg (a - b)]}, with the hint as on paper.
\end{itemize}
\end{frame}

%============================================================================
\section{Tactic proofs}

\begin{frame}[fragile]{Reading the infoview}
\begin{columns}[c]
\begin{column}{0.46\textwidth}
\begin{goal}
m n : ℕ
hm : Even m
hn : Even n
⊢ Even (m + n)
\end{goal}
\end{column}
\begin{column}{0.5\textwidth}
\begin{itemize}
\item Above the line: what you \textbf{have}. Variables, and hypotheses with names.
\item After \lc{⊢}: what you must \textbf{prove}, the \emph{goal}.
\item Each tactic changes this picture. The proof is finished when no goals are left.
\end{itemize}
\end{column}
\end{columns}

\vspace{1em}
\lean{even_def}
\begin{itemize}
\item In Mathlib, \lc{Even a} \emph{means} $\exists r,\ a = r + r$. \lc{#print} shows you the definition.
\end{itemize}
\end{frame}

%----------------------------------------------------------------------------
\begin{frame}[fragile,t]{A proof, step by step\livedemo}
\begin{columns}[T]
\begin{column}{0.5\textwidth}
\only<1>{\lean[highlightlines={1-3}, highlightcolor=goalbg]{even_proof}}%
\only<2>{\lean[highlightlines=4, highlightcolor=goalbg]{even_proof}}%
\only<3>{\lean[highlightlines=5, highlightcolor=goalbg]{even_proof}}%
\only<4>{\lean[highlightlines=6, highlightcolor=goalbg]{even_proof}}%
\only<5>{\lean[highlightlines=7, highlightcolor=goalbg]{even_proof}}%
\only<6>{\lean[highlightlines=8, highlightcolor=goalbg]{even_proof}}%

\vspace{0.4em}
\small
\only<1>{The starting position: two even numbers, and the goal.}%
\only<2>{\lc{obtain} unpacks \lc{hm : ∃ r, m = r + r}: we get a witness \lc{a} and the equation \lc{ha}.}%
\only<3>{The same for \lc{hn}.}%
\only<4>{To prove \lc{∃ r, m + n = r + r} we propose \lc{r := a + b}. (\lc{case h} is only a label.)}%
\only<5>{\lc{rw} rewrites \lc{m} and \lc{n} using the two equations.}%
\only<6>{\lc{ring} proves the identity. \textbf{No goals}: Lean accepts the proof.}%
\end{column}
\begin{column}{0.47\textwidth}
\begin{onlyenv}<1>
\begin{goal}
m n : ℕ
hm : Even m
hn : Even n
⊢ Even (m + n)
\end{goal}
\end{onlyenv}
\begin{onlyenv}<2>
\begin{goal}
m n : ℕ
hn : Even n
a : ℕ
ha : m = a + a
⊢ Even (m + n)
\end{goal}
\end{onlyenv}
\begin{onlyenv}<3>
\begin{goal}
m n a : ℕ
ha : m = a + a
b : ℕ
hb : n = b + b
⊢ Even (m + n)
\end{goal}
\end{onlyenv}
\begin{onlyenv}<4>
\begin{goal}
case h
m n a : ℕ
ha : m = a + a
b : ℕ
hb : n = b + b
⊢ m + n = a + b + (a + b)
\end{goal}
\end{onlyenv}
\begin{onlyenv}<5>
\begin{goal}
case h
m n a : ℕ
ha : m = a + a
b : ℕ
hb : n = b + b
⊢ a + a + (b + b) = a + b + (a + b)
\end{goal}
\end{onlyenv}
\begin{onlyenv}<6>
\begin{goal}
No goals
\end{goal}
\end{onlyenv}
\end{column}
\end{columns}
\end{frame}

%----------------------------------------------------------------------------
\begin{frame}[fragile]{The short way}
\lean{even_short}
\begin{itemize}
\item Mathlib already has this fact, under the name \lc{Even.add}. Writing \lc{hm.add hn} applies it to the two hypotheses.
\item Most short proofs you will see look like this: the work is in finding the right lemma.
\item More on finding lemmas in a moment.
\end{itemize}

\vspace{0.4em}
\begin{athome}
Type the step-by-step proof yourself, and watch the infoview after each line. Then \texttt{Exercises1.lean}, parts 1 and 2.
\end{athome}
\end{frame}

%----------------------------------------------------------------------------
\begin{frame}[fragile]{What Lean actually checks}
\begin{columns}[T]
\begin{column}{0.47\textwidth}
\lean{terms}
\begin{itemize}
\item A tactic is a program that \emph{builds} a proof term.
\item The kernel checks the term, not the tactic. That is why tactics, automation and AI can be trusted only as far as the kernel.
\end{itemize}
\end{column}
\begin{column}{0.5\textwidth}
\begin{leanout}[listing options={style=bgslean, basicstyle=\ttfamily\fontsize{3.3}{3.8}\selectfont}]
theorem small : ∀ (x : ℝ), 0 < x → 0 < x + 1 :=
fun x h =>
  lt_of_not_ge fun a =>
    Mathlib.Tactic.Linarith.lt_irrefl
      (Eq.mp
        (congrArg (fun _a => _a < 0)
          (Mathlib.Tactic.Ring.of_eq
            (Mathlib.Tactic.Ring.Common.add_congr
              (Mathlib.Tactic.Ring.Common.add_congr
                (Mathlib.Tactic.Ring.Common.neg_congr
                  (Mathlib.Tactic.Ring.cast_pos (Mathlib.Meta.NormNum.isNat_ofNat ℝ Nat.cast_one))
                  (Mathlib.Tactic.Ring.Common.neg_add
                    (Mathlib.Meta.NormNum.IsInt.to_raw_eq
                      (Mathlib.Meta.NormNum.isInt_neg (Eq.refl Neg.neg)
                        (Mathlib.Meta.NormNum.IsNat.to_isInt (Mathlib.Meta.NormNum.IsNat.of_raw ℝ 1))
                        (Eq.refl (Int.negOfNat 1))))
                    Mathlib.Tactic.Ring.Common.neg_zero))
                (Mathlib.Tactic.Ring.Common.sub_congr
                  (Mathlib.Tactic.Ring.cast_zero (Mathlib.Meta.NormNum.isNat_ofNat ℝ Nat.cast_zero))
                  (Mathlib.Tactic.Ring.Common.atom_pf x rfl
                    (Eq.mpr
                      (id
                        (congrArg (fun _a => x ^ Nat.rawCast 1 * Nat.rawCast 1 = x ^ Nat.rawCast 1 * _a) (Eq.symm rfl)))
                      (Eq.refl (x ^ Nat.rawCast 1 * Nat.rawCast 1))))
  ⋮   (68 lines in all)
\end{leanout}
\end{column}
\end{columns}
\end{frame}

%============================================================================
\section{Finding things in Mathlib}

\begin{frame}[fragile]{Ask Lean}
\begin{columns}[T]
\begin{column}{0.5\textwidth}
\lean{exact_q}
\begin{leanout}
Try this:
  [apply] exact two_mul_le_add_sq a b
\end{leanout}
\end{column}
\begin{column}{0.46\textwidth}
\begin{itemize}
\item \lc{exact?} searches Mathlib for a lemma that closes the goal.
\item \lc{apply?} finds lemmas whose conclusion fits; \lc{rw?} finds rewrites.
\item Click on the suggestion and it replaces the search.
\end{itemize}
\end{column}
\end{columns}

\vspace{0.5em}
\begin{caveat}[Before you prove it, search for it]
Our \lc{sq_add_sq_ge} from a few slides ago was already in Mathlib, as \lc{two_mul_le_add_sq}.
\end{caveat}
\end{frame}

%----------------------------------------------------------------------------
\begin{frame}[fragile]{When the search fails\livedemo}
\begin{columns}[T]
\begin{column}{0.48\textwidth}
\begin{leanout}[title={\scriptsize\sffamily Without a hypothesis}]
example (x : ℝ) :
    Real.sqrt x ^ 2 = x := by
  exact?
-- error: could not close the goal
\end{leanout}
\end{column}
\begin{column}{0.48\textwidth}
\lean{sqrt_fixed}
\begin{leanout}
Try this:
  [apply] exact Real.sq_sqrt hx
\end{leanout}
\end{column}
\end{columns}

\vspace{0.5em}
\begin{itemize}
\item The first statement is false: in Lean, $\sqrt{x} = 0$ for $x < 0$.
\item A search that fails is information. Before searching harder, ask whether the statement is true.
\end{itemize}
\end{frame}

%----------------------------------------------------------------------------
\begin{frame}[fragile]{Search engines for Mathlib}
\renewcommand{\arraystretch}{1.45}
\begin{tabular}{@{}>{\bfseries}p{0.15\linewidth}>{\raggedright\arraybackslash}p{0.4\linewidth}>{\raggedright\arraybackslash}p{0.37\linewidth}@{}}
\toprule
Tool & What you give it & Example\\\midrule
\lc{exact?} & a goal in your file & (previous slide)\\
Loogle & a pattern, with names and \lc{_} for holes & \lc{Real.sqrt, _ * _ = _}\\
LeanSearch & a sentence in English & \emph{square root of a product}\\
Mathlib docs & a name or a word & \lc{two_mul_le_add_sq}\\
\bottomrule
\end{tabular}

\vspace{0.8em}
\begin{itemize}
\item \href{https://loogle.lean-lang.org}{loogle.lean-lang.org} \quad \href{https://leansearch.net}{leansearch.net} \quad \href{https://leanprover-community.github.io/mathlib4_docs}{Mathlib documentation}
\item In the editor: hover over a name to see its statement; \emph{Go to definition} opens its source.
\end{itemize}
\end{frame}

%----------------------------------------------------------------------------
\begin{frame}[fragile]{Names follow the statement}
\lean{names}

\vspace{0.3em}
\begin{itemize}
\item A name reads the statement aloud: \lc{mul_comm} is about \lc{*} being commutative, \lc{mul_pos} about a product being positive.
\item For implications, the conclusion comes first: \lc{lt_of_le_of_lt} proves \lc{<} \emph{of} \lc{≤} and \lc{<}.
\item Dot notation: if \lc{hp : p.Prime}, then \lc{hp.two_le : 2 ≤ p}.
\item So you can often \emph{guess} a name, and check it with \lc{#check}.
\end{itemize}
\end{frame}

%----------------------------------------------------------------------------
\begin{frame}[fragile]{The library also has big theorems}
\lean{library}

\vspace{0.6em}
\begin{itemize}
\item Irrationality of $\sqrt2$, infinitely many primes, the fundamental theorem of algebra, \dots
\item Each of these is a name you can use in your own proofs.
\end{itemize}
\end{frame}

%============================================================================
\section{Reading a Lean statement}

\begin{frame}[fragile]{Three kinds of brackets}
\lean{signature}

\vspace{0.4em}
\renewcommand{\arraystretch}{1.3}
\begin{tabular}{@{}p{0.22\linewidth}>{\raggedright\arraybackslash}p{0.74\linewidth}@{}}
\lc{(a b : G)} & \textbf{explicit}: you supply them, \lc{mul_comm x y}.\\
\lc|{G : Type u_1}| & \textbf{implicit}: Lean works them out from the explicit ones.\\
\lc{[CommMagma G]} & \textbf{instance}: a structure on \lc{G}, found automatically. Here: a commutative multiplication.\\
\end{tabular}

\vspace{0.6em}
So \lc{mul_comm} applies to $\mathbb{N}$, $\mathbb{Z}$, $\mathbb{R}$, polynomial rings, \dots\ but not to square matrices: Lean finds no \lc{CommMagma} instance for them, and \lc{mul_comm A B} is rejected.
\end{frame}

%----------------------------------------------------------------------------
\begin{frame}[fragile]{Do not panic at long signatures}
\begin{leanout}
@two_mul_le_add_sq : ∀ {R : Type u_1} [inst : CommSemiring R]
  [inst_1 : LinearOrder R] [ExistsAddOfLE R] [MulPosStrictMono R]
  [AddLeftReflectLE R] [AddLeftMono R] (a b : R),
  2 * a * b ≤ a ^ 2 + b ^ 2
\end{leanout}

\vspace{0.5em}
\begin{itemize}
\item Read it as: \emph{for any ordered commutative semiring $R$, and $a, b \in R$, \ $2ab \le a^2 + b^2$.}
\item The instance arguments are the precise list of what the proof needs. $\mathbb{R}$ satisfies all of them.
\item Mathlib states results in this generality so that they apply to $\mathbb{N}$, $\mathbb{Q}$, $\mathbb{R}$, \dots at once.
\end{itemize}
\end{frame}

%----------------------------------------------------------------------------
\begin{frame}[fragile]{Junk values}
\begin{columns}[T]
\begin{column}{0.55\textwidth}
\lean{junk_basic}
\end{column}
\begin{column}{0.41\textwidth}
\begin{itemize}
\item In Lean every function is defined everywhere. Where mathematics says ``undefined'', Lean returns a convenient value, usually $0$.
\item Subtraction on $\mathbb{N}$ stops at $0$; division on $\mathbb{N}$ rounds down.
\end{itemize}
\end{column}
\end{columns}

\vspace{0.4em}
\begin{keypoint}
\normalsize This is harmless in theorems: the hypotheses do the work. A lemma about $a/b$ will ask for \lc{b ≠ 0}. It matters when \emph{you} write a statement.
\end{keypoint}
\end{frame}

%----------------------------------------------------------------------------
\begin{frame}[fragile]{Junk values in analysis}
\lean{junk_analysis}

\begin{itemize}
\item A theorem about $\sum' f$ that forgets \lc{Summable f} may hold only because the sum is $0$.
\item ``The least $n$ with property $B$'' is $0$ when no such $n$ exists.
\item ``$f$ has minimal period $N$'', written as ``no smaller period'', fails at $0$: every $f$ has period $0$.
\end{itemize}
\end{frame}

%----------------------------------------------------------------------------
\begin{frame}[fragile]{Coercions}
\lean{coercion}

\vspace{0.4em}
\begin{itemize}
\item \lc{((7 : ℕ) : ℝ)} is the natural number $7$ seen as a real number. Lean prints it with an arrow: \lc{↑n}.
\item Where the arrow sits changes the meaning: divide in $\mathbb{R}$ and you get $3.5$; divide in $\mathbb{N}$ first and you get $3$.
\item When a statement mixes $\mathbb{N}$, $\mathbb{Z}$ and $\mathbb{R}$, check where the arrows are.
\end{itemize}
\end{frame}

%----------------------------------------------------------------------------
\begin{frame}[fragile,t]{Spot the problem (1)}
\begin{paper}[Claim]
For all real $a, b$, \ $\dfrac{a}{b} \cdot b = a$.
\end{paper}

\pause
\lean{bug_div}

\begin{itemize}
\item The claim as written is \textbf{false} in Lean, and Lean can prove that it is false.
\item On paper the hypothesis $b \ne 0$ is often left implicit. In Lean it must be written.
\end{itemize}
\end{frame}

%----------------------------------------------------------------------------
\begin{frame}[fragile,t]{Spot the problem (2)}
\begin{paper}[Claim]
For every natural number $n$, \ $(n - 1) + 1 = n$.
\end{paper}

\pause
\lean{bug_nat}

\begin{itemize}
\item On $\mathbb{N}$, \ $0 - 1 = 0$. The claim fails at $n = 0$.
\item The same statement over $\mathbb{Z}$ is true. Which one did the paper mean?
\end{itemize}
\end{frame}

%----------------------------------------------------------------------------
\begin{frame}[fragile,t]{Spot the problem (3)}
\lean{vacuous}

\pause
\begin{caveat}[Lean accepts this]
The hypotheses $0 < x$ and $x < 0$ contradict each other, so \emph{anything} follows from them. The theorem is true and says nothing.
\end{caveat}

\begin{itemize}
\item A theorem with inconsistent hypotheses is proved for free. AI provers find such proofs quickly.
\item Check: can the hypotheses actually hold together? Try to give an example.
\end{itemize}
\end{frame}

%----------------------------------------------------------------------------
\begin{frame}[fragile,t]{Spot the problem (4)}
\lean{quantifiers}

\pause
\begin{itemize}
\item The only difference is where \lc{∀ x} sits.
\item \lc{Cont f} is continuity: $\delta$ may depend on $x$. \ \lc{UnifCont f} is uniform continuity: one $\delta$ for all $x$.
\item $x \mapsto x^2$ satisfies the first and not the second.
\end{itemize}

\begin{keypoint}
\normalsize Quantifier order is exactly the kind of detail a quick reading misses.
\end{keypoint}
\end{frame}

%----------------------------------------------------------------------------
\begin{frame}[fragile,t]{Quiz: which one is Goldbach's conjecture?}
\lean{goldbach_defs}

\only<1>{\vspace{0.5em}\begin{paper}[Goldbach's conjecture]
Every even integer greater than $2$ is the sum of two primes.
\end{paper}}%
\only<2>{\vspace{0.3em}\textbf{\color{alertred}Goldbach₂} forgets $n > 2$, and is false:
\lean{goldbach_two}}%
\only<3>{\vspace{0.3em}\textbf{\color{alertred}Goldbach₃} is true for a silly reason: \lc{∨} binds less tightly than \lc{∧}, so it says ``$p$ is prime, or ($q$ is prime and $n = p + q$)''.
\lean{goldbach_three}}%
\only<4>{\vspace{0.5em}\textbf{\color{okgreen}Goldbach₁} is the conjecture. Nobody knows a proof, so Lean cannot check one. But we can \emph{state} it, and check that the statement is right.}%
\end{frame}

%----------------------------------------------------------------------------
\begin{frame}[fragile]{It happens in practice}
From B.\ Alexeev's report on formalizing Erdős problems with AI:

\vspace{0.4em}
\begin{paper}[Problem 480]
\itshape ``The formalization of Problem 480 said \texttt{m ≠ 0} but meant to say \texttt{n ≠ 0}.''\ \upshape The AI found a counterexample that used the typo, together with a junk value.
\end{paper}
\begin{paper}[Problem 56]
\itshape ``Aristotle found a completely trivial counterexample with tiny parameters like 2. It turns out that the conjecture was missing a hypothesis, somewhat implicitly.''
\end{paper}

\begin{keypoint}
\normalsize ``Misformalization is a big issue and was more frequent than I had expected.''
\end{keypoint}

\source{B.\ Alexeev, guest post on K.\ Buzzard's Xena blog, 5 December 2025.}
\end{frame}

%----------------------------------------------------------------------------
\begin{frame}[fragile]{Trust, but check: \texttt{sorry} and axioms}
\lean{axioms}

\begin{itemize}
\item \lc{sorry} closes any goal and leaves a warning. A proof containing it is not a proof.
\item \lc{#print axioms} lists everything a theorem ultimately relies on.
\item The three standard axioms (\lc{propext}, \lc{Classical.choice}, \lc{Quot.sound}) are fine. \lc{sorryAx}, or an axiom you have never heard of, is not.
\end{itemize}

\begin{keypoint}
\normalsize For anything an AI hands you: check the statement, then \lc{#print axioms}.
\end{keypoint}
\end{frame}

%----------------------------------------------------------------------------
\begin{frame}[fragile]{Reading a statement: a checklist}
\begin{enumerate}
\item \textbf{Types.} $\mathbb{N}$, $\mathbb{Z}$ or $\mathbb{R}$? Where are the coercion arrows?
\item \textbf{Junk values.} Is there a division, a subtraction on $\mathbb{N}$, a square root, a logarithm, an infinite sum, an infimum? Is the case where it is ``undefined'' excluded?
\item \textbf{Hypotheses.} Are they all needed? Can they hold together?
\item \textbf{Quantifiers.} In which order? What may depend on what?
\item \textbf{Precedence.} How do \lc{∧}, \lc{∨}, \lc{→} group? Add brackets if unsure.
\item \textbf{Definitions.} Does each name mean what you think? Use \lc{#print} and hover.
\item \textbf{Proof.} No \lc{sorry}, only standard axioms.
\end{enumerate}
\end{frame}

%============================================================================
\section{At home}

\begin{frame}[fragile]{At home}
\begin{columns}[T]
\begin{column}{0.55\textwidth}
\texttt{Exercises1.lean}, in the browser or on your laptop.
\begin{enumerate}
\item Logic: \lc{intro}, \lc{use}, \lc{obtain}, \dots
\item Automation: pick the tactic from the table.
\item Search: find the Mathlib lemma.
\item Reading: true, false, or true for a silly reason? Decide, then prove it or its negation.
\end{enumerate}
\end{column}
\begin{column}{0.41\textwidth}
\begin{itemize}
\item First redo the three demos. Their code is in \texttt{Session1.lean}.
\item The tactic tables are at the end of these slides.
\item Stuck? Read the infoview first, then ask: on the Lean Zulip, or at the start of next session.
\item Want more? The \href{https://adam.math.hhu.de/\#/g/leanprover-community/nng4}{Natural Number Game}.
\end{itemize}
\end{column}
\end{columns}

\vspace{0.3em}
\begin{keypoint}
\normalsize Each exercise is rated twice: how hard it is to \emph{state}, and how hard to \emph{prove}. With AI, the first number is the one that matters.
\end{keypoint}
\end{frame}

%----------------------------------------------------------------------------
\begin{frame}[fragile]{Next time: from paper to Lean}
\begin{itemize}
\item Writing your own definitions and statements.
\item The usual mistakes, and how to catch them.
\item Getting an AI to draft a statement, and checking what it produces.
\end{itemize}

\vspace{1em}
\begin{keypoint}
\normalsize Before next time: think of a result from your own work you would like to see in Lean. Short and elementary is best: in session 4 I formalize some of them live.
\end{keypoint}
\end{frame}

%----------------------------------------------------------------------------
\begin{frame}[fragile]{Resources}
\begin{itemize}
\item J.\ Avigad and P.\ Massot, \emph{Mathematics in Lean}. The standard introduction for mathematicians.
\item J.\ Avigad et al., \emph{Theorem Proving in Lean 4}. How Lean works underneath.
\item The Natural Number Game. Lean as a puzzle game, in the browser.
\item The Mathlib documentation, Loogle and LeanSearch.
\item The Lean Zulip chat, stream \emph{\#new members}: questions of every level are welcome.
\end{itemize}
\end{frame}

%============================================================================
\appendix
\section{For the exercises}

%----------------------------------------------------------------------------
\begin{frame}[fragile]{Logic: the tactics you need}
\renewcommand{\arraystretch}{1.25}
\small
\begin{tabular}{@{}p{0.12\linewidth}>{\raggedright\arraybackslash}p{0.41\linewidth}>{\raggedright\arraybackslash}p{0.41\linewidth}@{}}
\toprule
& \textbf{to prove it} (in the goal) & \textbf{to use it} (a hypothesis \lc{h}) \\\midrule
\lc{P → Q}   & \lc{intro hp} & \lc{apply h}, or \lc{h hp}\\
\lc{∀ x, P x} & \lc{intro x} & \lc{h a}, or \lc{specialize h a}\\
\lc{∃ x, P x} & \lc{use a} & \lc{obtain ⟨x, hx⟩ := h}\\
\lc{P ∧ Q}   & \lc{constructor}, or \lc{exact ⟨hp, hq⟩} & \lc{obtain ⟨hp, hq⟩ := h}\\
\lc{P ∨ Q}   & \lc{left} or \lc{right} & \lc{rcases h with hp | hq}\\
\lc{¬ P}     & \lc{intro hp} (then prove \lc{False}) & \lc{exact absurd hp h}, \lc{contradiction}\\
\lc{P ↔ Q}   & \lc{constructor} & \lc{rw [h]}, or \lc{h.1} and \lc{h.2}\\
\lc{a = b}   & \lc{rfl}, \lc{ring}, \lc{simp}, \dots & \lc{rw [h]}\\
\bottomrule
\end{tabular}

\vspace{0.6em}
\normalsize
And always: \lc{exact h} when \lc{h} is exactly the goal.
\end{frame}

%----------------------------------------------------------------------------
\begin{frame}[fragile]{Logic: three small examples}
\begin{columns}[T]
\begin{column}{0.48\textwidth}
\lean{forall_exists}
\lean{and}
\end{column}
\begin{column}{0.48\textwidth}
\lean{or}
\end{column}
\end{columns}

\vspace{0.6em}
\begin{itemize}
\item \lc{omega} decides linear arithmetic over $\mathbb{N}$ and $\mathbb{Z}$; here it proves \lc{n < n + 1}.
\item \lc{⟨hq, hp⟩} builds a proof of a conjunction from its two halves.
\item \lc{·} focuses on one goal at a time. Indentation matters.
\end{itemize}
\end{frame}

%----------------------------------------------------------------------------
\begin{frame}[fragile]{Automation}
\lean{auto}

\vspace{0.2em}
\small
\begin{tabular}{@{}>{\ttfamily}p{0.11\linewidth}>{\raggedright\arraybackslash}p{0.34\linewidth}>{\ttfamily}p{0.11\linewidth}>{\raggedright\arraybackslash}p{0.34\linewidth}@{}}
ring & identities in commutative rings &
omega & linear arithmetic over $\mathbb{N}$, $\mathbb{Z}$\\
linarith & linear (in)equalities from hypotheses &
fun\_prop & continuity, measurability, \dots\\
positivity & goals \lc{0 < e}, \lc{0 ≤ e} &
gcongr & ``apply the same operation to both sides''\\
norm\_num & numerical facts &
simp & simplify using a large set of rewrite rules\\
\end{tabular}
\end{frame}

%----------------------------------------------------------------------------
\begin{frame}[fragile]{Propositions are types}
\renewcommand{\arraystretch}{1.35}
\begin{tabular}{@{}>{\raggedright\arraybackslash}p{0.27\linewidth}>{\raggedright\arraybackslash}p{0.33\linewidth}>{\raggedright\arraybackslash}p{0.36\linewidth}@{}}
\toprule
\textbf{Mathematics} & \textbf{Lean} & \textbf{Read as}\\\midrule
a statement $P$ & \lc{P : Prop} & $P$ is a proposition\\
a proof of $P$ & \lc{h : P} & $h$ is a proof of $P$\\
$P \Rightarrow Q$ & \lc{P → Q} & a function from proofs of $P$ to proofs of $Q$\\
for all $x \in \mathbb{R}$, $P(x)$ & \lc{∀ x : ℝ, P x} & a function from $x$ to proofs of $P(x)$\\
the theorem is true & \lc{theorem t : P := proof} & Lean accepted a term of type $P$\\
\bottomrule
\end{tabular}

\vspace{0.8em}
You will rarely need this explicitly. It explains why hypotheses are written \lc{(h : x < 0)}: a hypothesis is just a named proof you are given.
\end{frame}

\end{document}
